\subsection{Gelfand transformation}

Let A be a commutative Banach algebra. Recall that, for every \(x \in A\) , we denote by \(\mathcal{G}_{\mathrm{A}}(x)\) , or \(\mathcal{G}(x)\) , the function \(\chi \mapsto \chi(x)\) on \(\mathsf{X}(\mathsf{A})\) , that \(\mathcal{G}(x)\) is called the Gelfand transform of x, and that the map \(x \mapsto \mathcal{G}(x)\) is called the Gelfand transformation (cf.~def. 5 of I, p.~7). One therefore has by definition:

\[
\mathcal {G} (x) (\chi) = \chi (x).
\]

Examples. --- 1) Let X be a locally compact topological space and consider the commutative Banach algebra \(\mathcal{C}_{0}(\mathrm{X})\) (example 3 of I, p.~17 and number 2 of I, p.~31). By cor. 1 of I, p.~32, the space of characters \(\mathsf{X}(\mathcal{C}_{0}(\mathbf{X}))\) is identified with X by means of the map associating to an element \(x \in \mathbf{X}\) the character \(f \mapsto f(x)\) of \(\mathcal{C}_{0}(\mathbf{X})\) , and the Gelfand transformation of \(\mathcal{C}_{0}(\mathbf{X})\) is then identified with the identity map.

\begin{enumerate}
\def\labelenumi{\arabic{enumi})}
\setcounter{enumi}{1}
\tightlist
\item
  Let \(n \geqslant 0\) be an integer. Let \(A_n\) be the algebra of functions \(f: [0,1] \to K\) admitting continuous derivatives in \([0,1]\) up to order \(n\) . Equipped with the norm
\end{enumerate}

\[
\| f \| = \sum_ {k = 0} ^ {n} \frac {1}{k !} \sup _ {0 \leqslant t \leqslant 1} | f ^ {(k)} (t) |,
\]

it is a Banach algebra (example 4 of I, p.~18). For every n, the space of characters \(\mathsf{X}(\mathsf{A}_{n})\) is identified with {[}0,1{]} and G under inclusion \(A_{n}\) into \(\mathcal{C}([0,1])\) (cf.~example 1 of I, p.~144).

\begin{enumerate}
\def\labelenumi{\arabic{enumi})}
\setcounter{enumi}{2}
\item
  Let \(\Delta\) be the disk of complex numbers \(z\) satisfying \(|z| \leqslant 1\) and let A be the complex Banach algebra of continuous functions on \(\Delta\) that are analytic in the interior of \(\Delta\) , equipped with the norm \(\|f\| = \sup_{z \in \Delta} |f(z)|\) (example 9 of I, p.~20). Then \(\mathsf{X}(\mathsf{A})\) is identified with \(\Delta\) and \(\mathcal{G}\) with the inclusion of A into \(\mathcal{C}(\Delta)\) (cf.~exerc. 6 of I, p.~193).
\item
  Consider the complex Banach algebra A of absolutely convergent Fourier series (example 8 of I, p.~19). For every element \(u\) of the unit circle U, the map \(f \mapsto f(u)\) is a character \(ev_{u}\) of A. If \(f_{0} \in A\) is the identity map of U, one has \(ev_{u}(f_{0}) = u\) , hence the map ev: \(u \mapsto ev_{u}\) from U into X(A) is injective; it is continuous. Let \(\chi \in \mathsf{X}(\mathsf{A})\) . We have \(\|f_{0}\| = \|f_{0}^{-1}\| = 1\) , hence \(|\chi(f_{0})| \leqslant 1\) and \(|\chi(f_{0})^{-1}| \leqslant 1\) . This shows that \(\chi(f_{0}) \in \mathbf{U}\) and there exists \(u \in U\) such that \(\chi(f_{0}) = \mathrm{ev}_{u}(f_{0})\) . Since \(\{f_{0}, f_{0}^{-1}\}\) topologically generates the unital algebra A, one has \(\chi = ev_{u}\) . Thus, the map ev is a homeomorphism from U onto X(A), by which one identifies these spaces. The Gelfand transformation of A is then identified with the inclusion of A into \(\mathcal{C}(\mathbf{U})\) .
\end{enumerate}

Since A is isomorphic to the Banach algebra \(\mathrm{L}^{1}(\mathbf{Z})\) (example 8 of I, p.~19), the space \(\mathsf{X}(\mathrm{L}^{1}(\mathbf{Z}))\) is identified with U and, for every element \((c_{n})\in\mathrm{L}^{1}(\mathbf{Z})\) , the Gelfand transform \(\mathcal{G}_{\mathrm{L}^{1}(\mathbf{Z})}((c_{n}))\) is identified with the function \(u\mapsto\sum_{n\in\mathbf{Z}}c_{n}u^{n}\) on U.

\begin{enumerate}
\def\labelenumi{\arabic{enumi})}
\setcounter{enumi}{4}
\tightlist
\item
  Let \(\Delta\) the unit disk of complex numbers \(z\) such that \(|z| \leqslant 1\) . Its boundary in \(\mathbf{C}\) is \(\mathbf{U}\) . Let A be the Banach algebra of complex functions \(f\) on \(\mathbf{U}\) such that there exists a continuous function \(\widetilde{f} \in \mathcal{C}(\Delta)\) extending \(f\) that is analytic in \(\mathring{\Delta}\) , equipped with the norm \(\|f\| = \sup_{z \in \mathbf{U}} |f(z)|\) . By virtue of the maximum principle (VAR, R1, p.~30, 3.3.7), one then has \(\|f\| = \sup_{z \in \Delta} |\widetilde{f}(z)|\) , and hence A coincides with the algebra of example 9 of I, p.~20. The set \(\mathsf{X}(\mathsf{A})\) is identified with \(\Delta\) and, if \(f \in \mathsf{A}\) , the map \(\mathcal{G}(f)\) is identified with the continuous extension of \(f\) into \(\Delta\) that is analytic in \(\mathring{\Delta}\) (cf.~exerc. 6 of I, p.~193).
\end{enumerate}

PROPOSITION 5. --- Let A be a commutative Banach algebra. For every \(x \in \mathbf{A}\) , the function \(\mathcal{G}(x)\) belongs to the commutative Banach algebra \(\mathcal{C}_0(\mathsf{X}(\mathbf{A}))\) of continuous functions on \(\mathsf{X}(\mathbf{A})\) vanishing at infinity.

By definition (cf.~no. 7 of I, p.~9), the function \(\mathcal{G}_{\mathrm{A}}^{\prime}(x)\colon \chi \mapsto \chi (x)\) is continuous on \(\mathsf{X}'(\mathsf{A})\) and vanishes at 0. Since \(\mathsf{X}'(\mathsf{A})\) is identified with the Alexandroff compactification of \(\mathsf{X}(\mathsf{A})\) according to cor. 1 of I, p.~29, the proposition follows.

PROPOSITION 6. --- Let A be a commutative Banach algebra and let \(x \in A\) .

\begin{enumerate}
\def\labelenumi{\alph{enumi})}
\item
  The union of the set of values of \(\mathcal{G}(x)\) and of \(\{0\}\) is equal to \(\mathrm{Sp}_{\mathrm{A}}^{\prime}(x)\) ;
\item
  If A has a unit element, the set of values of \(\mathcal{G}(x)\) is \(\mathrm{Sp}_{\mathrm{A}}(x)\) . In particular, for \(x\) to be invertible, it is necessary and sufficient that \(\mathcal{G}(x)\) does not vanish.
\end{enumerate}

Suppose that A has a unity element. We know that, for every \(\chi \in \mathsf{X}(\mathsf{A})\) , one has \(\chi(x) \in \mathrm{Sp}_{\mathsf{A}}(x)\) Conversely, let \(\lambda \in \mathrm{Sp}_{\mathsf{A}}(x)\) . Then \(x - \lambda\) is not invertible, hence belongs to a maximal ideal of A. There then exists \(\chi \in \mathsf{X}(\mathsf{A})\) such that \(\chi(x - \lambda) = 0\) (th. 2 of I, p.~30), whence \(b\) ).

Moving to the general case. Let \(\widetilde{\mathsf{A}}\) the Banach algebra obtained from A by adjoining an identity element; it is commutative. The set \(\mathrm{Sp}_{\widetilde{\mathsf{A}}}^{\prime}(x)\) is equal to \(\mathrm{Sp}_{\widetilde{\mathsf{A}}}^{\sim}(x)\) , that is, the set of values of \(\mathcal{G}_{\widetilde{\mathsf{A}}}^{\sim}(x)\) on \(\mathsf{X}(\widetilde{\mathsf{A}}) = \mathsf{X}'(\mathsf{A})\) Hence \(a\) ).

Example. --- Consider the Banach algebra A of absolutely convergent Fourier series (example 4). Prop. 6, b) implies that if \(\varphi\) is a function on the unit circle U admitting an absolutely convergent Fourier series, and if \(\varphi\) does not vanish, the function \(1/\varphi\) also admits an absolutely convergent Fourier series (“Wiener’s theorem”).

PROPOSITION 7. --- Let A be a commutative Banach algebra.

\begin{enumerate}
\def\labelenumi{\alph{enumi})}
\item
  The Gelfand transformation \(\mathcal{G}\) defines a morphism from A to \(\mathcal{C}_0(\mathsf{X}(\mathsf{A}))\) such that \(\| \mathcal{G}(x)\| = \varrho (x)\leqslant \| x\|\) for all \(x\in \mathrm{A}\) ;
\item
  For the Gelfand transform to \(\mathcal{G}\) it is isometric, it is necessary and sufficient that \(\|x^{2}\|=\|x\|^{2}\) for all \(x\in\mathrm{A}\) .
\end{enumerate}

The map \(\mathcal{G}\) is a morphism from A to \(\mathcal{C}_0(\mathsf{X}(\mathsf{A}))\) by no. 7 of I, p.~9 and prop. 5 of I, p.~37, and verifies \(\| \mathcal{G}(x)\| = \varrho (x)\) by prop. 6 and cor. 5 of I, p.~28. The assertion \(b)\) follows from \(a)\) and of Remark 1 of I, p.~21.

COROLLARY. --- Let A be a Banach algebra, and let x and y be commuting elements of A.

\begin{enumerate}
\def\labelenumi{\alph{enumi})}
\item
  We have \(\varrho(xy) \leqslant \varrho(x)\varrho(y)\) and \(\varrho(x + y) \leqslant \varrho(x) + \varrho(y)\) ;
\item
  If y is quasi-nilpotent, then \(\mathrm{Sp}_{\mathrm{A}}^{\prime}(x)=\mathrm{Sp}_{\mathrm{A}}^{\prime}(x+y)\) ; if moreover A is unital, then \(\mathrm{Sp}_{\mathrm{A}}(x)=\mathrm{Sp}_{\mathrm{A}}(x+y)\) .
\end{enumerate}

By considering the Banach algebra obtained from A by adjoining an identity element, one first reduces to the case where the algebra A is unital. Then, by considering the closed full subalgebra of A generated by x and y, one reduces to the case where A is commutative and unital. Assertion (a) is then a consequence of Prop. 7(a), and assertion (b) follows from Prop. 6 of I, p.~37 and Cor. 5 of I, p.~28.

PROPOSITION 8. --- Let A be a commutative Banach algebra. The following four sets are equal:

\begin{enumerate}
\def\labelenumi{(\roman{enumi})}
\item
  The kernel of the Gelfand transformation;
\item
  The set of elements x of A such that \(\mathrm{Sp}_{\mathrm{A}}^{\prime}(x)=\{0\}\) ;
\item
  The set of quasi-nilpotent elements of A;
\item
  The radical of A.
\end{enumerate}

We denote \(N_{1}\) , \(N_{2}\) , \(N_{3}\) , \(N_{4}\) , respectively, these sets. We have \(N_{1} = N_{2}\) (prop. 6, a)), and \(N_{2} = N_{3}\) (I, p.~28, cor. 5). By definition, the set \(N_{4}\) is the intersection of the regular maximal ideals of A; it is therefore the intersection of the kernels of the characters of A (th. 2 of I, p.~30), which is equal to \(N_{1}\) .

Remarks. --- 1) In general, the image of the Gelfand transformation is neither closed in \(\mathcal{C}_{0}(\mathsf{X}(\mathsf{A}))\) , nor dense in \(\mathcal{C}_{0}(\mathsf{X}(\mathsf{A}))\) (exerc. 7 of I, p.~193).

\begin{enumerate}
\def\labelenumi{\arabic{enumi})}
\setcounter{enumi}{1}
\item
  The image of the Gelfand transformation separates the points of \(\mathsf{X}(\mathsf{A})\) , since if \(\chi_1 \neq \chi_2\) are elements of \(\mathsf{X}(\mathsf{A})\) , there exists \(x \in \mathsf{A}\) such that \(\chi_1(x) \neq \chi_2(x)\) .
\item
  If \(\chi \in \mathsf{X}(\mathsf{A})\) , there exists an element of the image of the Gelfand transformation which does not vanish at \(\chi\) .
\item
  If A has an identity element, the image of the Gelfand transformation is a full subalgebra of the algebra of continuous functions on \(\mathsf{X}(\mathsf{A})\) (prop. 6, b)).
\end{enumerate}

Lemma 2. --- Let A be a commutative Banach algebra. Let M be a subset of X(A). Then M is closed for the Jacobson topology if and only if, for every character \(\chi \in \mathsf{X}(\mathsf{A}) = \mathsf{M}\) , there exists an element \(x\) of A such that \(\mathcal{G}(x)\) vanishes on M and is nonzero at \(\chi\) .

Let \(\Upsilon(\mathrm{M})\) the intersection of the kernels of the elements of M. The set M is closed for the Jacobson topology if and only if \(\mathrm{M} = \mathrm{V}(\Upsilon(\mathrm{M}))\) (cf.~I, p.~13 and I, p.~30). This condition is equivalent to saying that the elements \(\chi\) of M are precisely the characters which vanish on \(\Upsilon(\mathrm{M})\) . Consequently M is closed if and only if for every character \(\chi \notin \mathrm{M}\) , there exists \(x \in \Upsilon(\mathrm{M})\) such that \(\chi(x) \neq 0\) . This translates into \(\mathcal{G}(x)(\chi) \neq 0\) and \(\mathcal{G}(x)|_{\mathrm{M}} = 0\) .

